% =========================
% Explanation: Reproducibility (informal, readable)
% =========================
\subsection*{Explanation of Reproducibility Statement (Exact vs.\ Broad) and Artifact Plan (informal)}
This section explains which parts of the study can be reproduced exactly, which can only be reproduced within statistical tolerance,
and what artifacts we release so others can rerun the experiments and regenerate all reported results.

\begin{itemize}
    \item \textbf{Why distinguish ``exact'' vs.\ ``broad'' reproducibility?}
    In HPO/NAS, reproducibility depends on whether objective evaluations are deterministic.
    Tabular benchmarks or fully deterministic pipelines can often be reproduced \emph{bitwise}.
    GPU training often has unavoidable nondeterminism (e.g., non-deterministic kernels, floating-point non-associativity, mixed precision),
    so repeated runs with the same seed may differ slightly.
    We therefore label which regime applies to each experiment to avoid overstating reproducibility.

    \item \textbf{Exact reproducibility (bitwise-identical reruns).}
    Experiments are labeled \emph{exactly reproducible} when:
    \begin{itemize}
        \item the evaluation function is deterministic given \(\lambda\) and a seed (e.g., tabular lookup, fixed CV with deterministic learners),
        \item the software stack supports deterministic execution end-to-end (deterministic kernels, fixed dataloader order),
        \item and all sources of randomness are controlled and recorded.
    \end{itemize}
    In this regime, releasing code, configs, seeds, and environment pins should enable identical results.

    \item \textbf{Broad reproducibility (statistical reproducibility).}
    Experiments are labeled \emph{broadly reproducible} when objective evaluations are stochastic or hardware-dependent,
    so reruns may not be bitwise identical.
    Here, reproducibility means that:
    (i) rerunning the same protocol on the same stack yields results consistent within the reported variability, and
    (ii) the main conclusions (rankings, effect sizes, confidence intervals) remain stable.
    We therefore emphasize distributions across seeds and include repeated-seed re-evaluations of final configurations where appropriate.

    \item \textbf{What we release (artifacts required for reproduction).}
    To enable regeneration of all paper figures and tables, we release artifacts covering:
    \begin{itemize}
        \item \textbf{Code:} the proposed method and all runnable baselines (or wrappers around pinned upstream implementations).
        \item \textbf{Unified harness:} a shared evaluator/runner enforcing identical budgets, stopping rules, failure handling,
        seed policy, and parallelism semantics across methods.
        \item \textbf{Configuration files:} machine-readable configs for search spaces, datasets/splits, training recipes,
        method hyperparameters, budgets, and seed lists.
        \item \textbf{Environment specification:} Docker image / digest or Conda lockfile, plus explicit versions of key libraries
        (e.g., CUDA/cuDNN/PyTorch) and relevant hardware assumptions.
        \item \textbf{Data provenance:} dataset versions/snapshots, download scripts where permitted, and split manifests
        so that the same examples are used.
        \item \textbf{Raw trial logs:} per-evaluation records including \(\lambda\), objective values, timestamps/evaluation index,
        seeds, fidelity level (if applicable), and failure codes.
        \item \textbf{Aggregation and plotting scripts:} scripts that consume raw logs and regenerate all tables/figures,
        including statistical tests and multiple-comparisons corrections.
        \item \textbf{Run manifest:} for every run, record benchmark id, method, budget, seeds, worker count, git hash, container hash,
        and (when available) hardware identifier for provenance.
    \end{itemize}

    \item \textbf{Reproducibility checklist items (what we document).}
    We document all details needed to interpret and reproduce results:
    budgets, number of repetitions, seed policies, stopping rules,
    handling of invalid/failed trials, and any deviations from reference baseline implementations.
    We also document known sources of nondeterminism and quantify their impact (e.g., variance across reruns).

    \item \textbf{Comparability and limitations.}
    Not all methods are directly comparable under the same evaluation model.
    We state explicitly:
    \begin{itemize}
        \item which baselines are \emph{directly comparable} (same search space, same evaluation unit, same budget accounting),
        \item which are excluded or only reported under a separate setting (e.g., weight-sharing NAS, methods requiring different fidelities),
        \item and how these feasibility constraints bound claims (no claims about methods that are not directly comparable).
    \end{itemize}
    This prevents overreach and clarifies what the evidence supports.
\end{itemize}
